\documentclass[11pt]{article}

\usepackage{amsmath}
\usepackage{mathtools}
\usepackage{amsthm}
\usepackage[hidelinks]{hyperref}
\usepackage{cleveref}

\theoremstyle{plain}
\newtheorem{proposition}{Proposition}
\theoremstyle{definition}
\newtheorem{definition}{Definition}

\title{A Note on JOIN Multiplicity in SPARQL Conjunctive Queries}
\author{Bryan-Elliott Tam}
\date{June 4, 2026}

\begin{document}
\maketitle

\section{Introduction}

SPARQL is evaluated under bag-set semantics: an RDF knowledge graph (KG) is a set of triples, while SPARQL operators return \emph{bags} (multisets) of solution mappings.
A conjunctive query (CQ) is a query whose body is a single basic graph pattern (BGP), optionally followed by a projection; its only body operator is \texttt{JOIN}.

This note proves that, in a CQ, \texttt{JOIN} cannot raise the multiplicity of a solution mapping above one.
Within a CQ, projection at the head is therefore the only operator that can introduce duplicates.
The argument is SPARQL-specific: it works directly from the SPARQL specification's definition of \texttt{JOIN} multiplicity~\cite{w3SPARQLQuery} over solution mappings, and gives the SPARQL counterpart of the relational result of Chaudhuri and Vardi~\cite{Chaudhuri1993} and Afrati et al.~\cite{Afrati2010} that, for CQs, only projection increases multiplicity.

\section{Preliminaries}

\begin{definition}[Multiplicity~\cite{w3SPARQLQuery}]\label{def_multiplicity}
Given a bag (multiset) $\Omega$ of solution mappings and a solution mapping $\mu$, we write $\operatorname{multiplicity}(\mu \mid \Omega)$ for the number of times $\mu$ occurs in $\Omega$.
\end{definition}

\begin{definition}[Merge]\label{def_merge}
Two solution mappings $\mu_1$ and $\mu_2$ are \emph{compatible} if $\mu_1(v)$ and $\mu_2(v)$ are the same RDF term for every variable $v \in \operatorname{dom}(\mu_1) \cap \operatorname{dom}(\mu_2)$.
The \emph{merge} of two compatible solution mappings is their union, written $\operatorname{merge}(\mu_1, \mu_2) = \mu_1 \cup \mu_2$.
\end{definition}

\begin{definition}[SPARQL \texttt{JOIN} multiplicity~\cite{w3SPARQLQuery}]\label{def_join_multiplicity}
The multiplicity of a solution mapping $\mu$ in the SPARQL \texttt{JOIN} of two bags $\Omega_1$ and $\Omega_2$ is
\begin{multline*}
\operatorname{multiplicity}(\mu \mid \texttt{JOIN}(\Omega_1, \Omega_2)) = \\
\sum_{(\mu_1, \mu_2) \in m(\mu)} \operatorname{multiplicity}(\mu_1 \mid \Omega_1) \cdot \operatorname{multiplicity}(\mu_2 \mid \Omega_2),
\end{multline*}
where $m(\mu) = \{(\mu_1, \mu_2) \mid \operatorname{merge}(\mu_1 \in \Omega_1, \mu_2 \in \Omega_2) = \mu\}$.
\end{definition}

\section{Result}

\begin{proposition}\label{prop_cq_join_no_dup}
In a CQ, \texttt{JOIN} cannot produce a solution mapping with multiplicity higher than one.
\end{proposition}

\begin{proof}
Let $\Omega_1$ and $\Omega_2$ be the input bags of solution mappings of the \texttt{JOIN}.
Each input bag either comes directly from the set-semantic KG or is the output of an earlier \texttt{JOIN}.
By \Cref{def_join_multiplicity}, $\operatorname{multiplicity}(\mu \mid \texttt{JOIN}(\Omega_1, \Omega_2))$ exceeds one only if (i) the index set $m(\mu)$ contains more than one pair, or (ii) some product $\operatorname{multiplicity}(\mu_1 \mid \Omega_1) \cdot \operatorname{multiplicity}(\mu_2 \mid \Omega_2)$ exceeds one.
We rule out both in the context of a CQ.

\paragraph{Case (i).}
For $|m(\mu)| > 1$, two distinct pairs of solution mappings must merge to the same $\mu$, which by \Cref{def_merge} requires an input bag to contain solution mappings with different domains, for instance $\Omega_1 = [\{?x \mapsto a\}, \{?y \mapsto b, ?x \mapsto a\}]$ and $\Omega_2 = [\{?y \mapsto b, ?x \mapsto a\}, \{?x \mapsto a\}]$.
No CQ bag can have this shape: a triple pattern, or a BGP, binds a fixed set of variables, and a \texttt{JOIN} returns the merges of compatible solution mappings, each defined over the union of the two input domains.
Every bag therefore has a single fixed domain, so $|m(\mu)| \leq 1$.

\paragraph{Case (ii).}
The multiplicity of $\mu$ is the sum over $m(\mu)$ of the products $\operatorname{multiplicity}(\mu_1 \mid \Omega_1) \cdot \operatorname{multiplicity}(\mu_2 \mid \Omega_2)$.
Case (i) holds at every \texttt{JOIN} of the CQ, so no index set exceeds one and each such sum reduces to a single product.
Unfolding these products down to the KG leaves the multiplicity of $\mu$ as a product of solution mapping multiplicities taken directly from the set-semantic KG, each of which is one.
Hence the multiplicity of $\mu$ is one.
\end{proof}

\bibliographystyle{plain}
\bibliography{references}

\end{document}
